% =====================================================================
%  preview.tex -- standalone two-column preview of the accounting section.
%  Compile:  cd theory && pdflatex preview.tex   (twice, for refs)
%  The deliverable for the paper is section_privacy_accounting.tex, which
%  compiles inside IEEEtran; this file only proves it compiles at all and
%  shows the intended two-column layout with the PGFPlots figures.
% =====================================================================
\documentclass[10pt,twocolumn,letterpaper]{article}

\usepackage[T1]{fontenc}
\usepackage[utf8]{inputenc}
\usepackage{amsmath,amssymb,amsthm}
\usepackage{booktabs}
\usepackage{graphicx}
\usepackage{pgfplots}
\usepackage{pgfplotstable}
\usepackage[margin=0.75in]{geometry}
\usepackage[colorlinks=true,linkcolor=black,urlcolor=blue]{hyperref}

% ---- stand-ins for the IEEEtran-only bits used by the section ----------
\providecommand{\RDP}{\mathrm{RDP}}
\providecommand{\IEEEPARstart}[2]{#1#2\ }

\def\figdatadir{../figures/dat/}

% IEEE-style index terms environment (stand-in for IEEEtran's IEEEkeywords)
\newenvironment{IEEEkeywords}{\paragraph{Index Terms:}}{}

\begin{document}

\title{Semantic Communications with Differential Privacy: Dynamic Semantic
Budget Allocation and a Bookable Privacy Accounting}
\author{Q.~Wang \and A.~Coauthor}
\date{}
\maketitle

\begin{abstract}
Semantic communication links transmit latent codes that may carry identifying
information about the source subject. Existing privacy-oriented designs
justify protection by the marginal distribution of the released noise rather
than by an auditable privacy budget, which leaves the guarantee hard to state
and compare. This paper develops a bookable accounting for decomposed semantic
latents and an accompanying noise allocation. We isolate the per-component L2
sensitivity, show how the accounting rule fixes the optimal noise profile, and
give a closed-form water-filling KKT solution with two asymptotic regimes. A
Gaussian worst-case mutual-information bound combined with Fano's inequality
certifies an identification-error floor for an adversary observing the
release, and we derive the complete R\'enyi-DP to $(\varepsilon,\delta)$
conversion and the Gaussian-DP quadrature composition. On a synthetic
calibration anchored to the reference backbone, a $0.90$ certified floor over
$10^{4}$ templates leaves a $0.33$-bit joint leakage budget; quadrature
composition matches the additive-noise mechanism and is tighter than linear
composition, whose amplitude-only form under-reports $\varepsilon$ at high SNR
($693$ versus $2441$ at the audited release). We also show that an invertible
learned shaping cannot raise the guarantee: post-processing leaves the privacy
loss unchanged, so such a component is legitimate only as a keyed design
element.
\end{abstract}

\begin{IEEEkeywords}
Semantic communications, differential privacy, R\'enyi differential privacy,
Gaussian differential privacy, mutual information, Fano inequality, resource
allocation, water-filling, latent representations.
\end{IEEEkeywords}

% =====================================================================
%  figures_pgfplots.tex  --  camera-ready vector figures for the
%  privacy-accounting section.  Every plot reads a .dat file written by
%  `python -m sdpa.figures`, so text, ticks and fonts are LaTeX-native.
%
%  Set the data directory before \input, e.g.
%      \def\figdatadir{../figures/dat/}
%  The default below assumes you compile from inside figures/.
% =====================================================================
\providecommand{\figdatadir}{dat/}
\pgfplotsset{compat=1.18}

% ---------------------------------------------------------------------
% Fig. 1 -- the accounting law dictates the noise profile
% ---------------------------------------------------------------------
\begin{figure}[t]
\centering
\begin{tikzpicture}
  \pgfplotstableread{\figdatadir budget_split.dat}\bs
  \begin{axis}[
      width=\columnwidth, height=0.70\columnwidth,
      xlabel={semantic component $h$},
      ylabel={share of noise power},
      xmin=-0.8, xmax=14.4, ymin=0, ymax=0.30,
      xtick={0,2,...,14},
      tick label style={font=\footnotesize},
      label style={font=\footnotesize},
      legend style={font=\scriptsize, draw=none, fill=none,
                    at={(0.5,-0.26)}, anchor=north, legend columns=3},
      ymajorgrids, grid style={densely dotted, black!25},
      extra x ticks={0,1,3,4,5,6},
      extra x tick labels={},
      extra tick style={grid=major, major grid style={black!15}},
    ]
    \addplot[draw=none, fill=gray!10] coordinates {
      (1.5,0)(2.5,0)(2.5,0.30)(1.5,0.30)(1.5,0)
      (6.5,0)(7.5,0)(7.5,0.30)(6.5,0.30)(6.5,0)
      (7.5,0)(8.5,0)(8.5,0.30)(7.5,0.30)(7.5,0)
      (8.5,0)(9.5,0)(9.5,0.30)(8.5,0.30)(8.5,0)
      (9.5,0)(10.5,0)(10.5,0.30)(9.5,0.30)(9.5,0)
      (10.5,0)(11.5,0)(11.5,0.30)(10.5,0.30)(10.5,0)
      (11.5,0)(12.5,0)(12.5,0.30)(11.5,0.30)(11.5,0)
      (12.5,0)(13.5,0)(13.5,0.30)(12.5,0.30)(12.5,0)
      (13.5,0)(14.5,0)(14.5,0.30)(13.5,0.30)(13.5,0)};
    \addplot+[ybar, bar width=4pt, fill=blue!55!black, draw=none]
      table[x expr=\thisrow{h}-0.27, y=gdp] {\bs};
    \addplot+[ybar, bar width=4pt, fill=blue!30, draw=none]
      table[x expr=\thisrow{h}-0.09, y=rdp] {\bs};
    \addplot+[ybar, bar width=4pt, fill=orange!70!black, draw=none]
      table[x expr=\thisrow{h}+0.09, y=basic] {\bs};
    \addplot+[ybar, bar width=4pt, fill=black!40, draw=none]
      table[x expr=\thisrow{h}+0.27, y=static] {\bs};
    \addplot[black, mark=*, mark size=1pt, densely dotted, mark options={solid}]
      table[x=h, y=wkappa_norm] {\bs};
    \legend{GDP, RDP, $\ell_1$ (ampl.), static ($k{=}4$), $w_h\kappa_h$}
  \end{axis}
\end{tikzpicture}
\caption{Noise-power share of the four accounting rules on the task-decoupled
calibration ($\varepsilon{=}100$, $\delta{=}10^{-5}$; this panel illustrates
the shape of the allocation, while the audit in Fig.~\ref{fig:composite}(b)
uses the $\varepsilon{=}2000$ operating point). Only the declared private
components $h\in P=\{0,1,3,4,5,6\}$ may carry noise
(Theorem~\ref{th:allocation}); the gray-shaded blocks are the task complement
$Q$. The amplitude-only $\ell_1$ rule concentrates noise on a few high-sensitivity
components, the quadrature GDP/RDP rules spread it more evenly, and the static
rule adds noise to only the four selected blocks and leaves the rest untouched
(Corollary~\ref{co:static}).}
\label{fig:budget_split}
\end{figure}

% ---------------------------------------------------------------------
% Fig. 2 -- the price of a certificate
% ---------------------------------------------------------------------
\begin{figure}[t]
\centering
\begin{tikzpicture}
  \pgfplotstableread{\figdatadir certificate_price.dat}\cp
  \begin{axis}[
      name=left, width=0.84\columnwidth, height=0.50\columnwidth,
      xlabel={certified error floor $\tau$},
      ylabel={Bob LPIPS},
      xmin=0.29, xmax=0.96, ymin=0.10, ymax=1.0,
      tick label style={font=\footnotesize},
      label style={font=\footnotesize},
      legend style={font=\scriptsize, draw=none, fill=none,
                    at={(0.5,-0.24)}, anchor=north, legend columns=2,
                    nodes={font=\scriptsize, inner sep=1pt}},
      ymajorgrids, grid style={densely dotted, black!25},
    ]
    \addplot[red!70!black, mark=square*, mark size=1.2pt]
      table[x=tau, y=lpips_aligned] {\cp};
    \addplot[blue!70!black, mark=*, mark size=1.2pt]
      table[x=tau, y=lpips_decoupled] {\cp};
    \addplot[black, densely dotted, domain=0.29:0.96] {0.112};
    \legend{{Bob, aligned $w_h$}, {Bob, decoupled $w_h$},
            {no-protection $0.112$}}
  \end{axis}
  \begin{axis}[
      name=right, at={(left.south west)}, anchor=south west,
      width=0.84\columnwidth, height=0.50\columnwidth,
      axis y line*=right, axis x line=none, xtick=\empty,
      ylabel={identification error $P_e$},
      ymin=0.25, ymax=1.15,
      tick label style={font=\footnotesize},
      label style={font=\footnotesize},
    ]
    \addplot[blue!70!black, densely dotted, very thick]
      table[x=tau, y=certified_floor] {\cp};
    \addplot[blue!70!black, only marks, mark=triangle*, mark size=1.4pt,
             densely dashed]
      table[x=tau, y=p_e_surrogate] {\cp};
  \end{axis}
\end{tikzpicture}
\caption{Certificate price on the synthetic anchor, two overlap regimes. The
two Bob LPIPS curves are read on the \emph{left} axis; the dotted
certified-floor curve and the triangle surrogate $P_e$ are read on the
\emph{right} axis. The LPIPS values are surrogate predictions from the default
distortion model, not measured reconstruction distortion. When the private
semantics \emph{are} the task (aligned, $w_h$ large, task overlap $\approx1$)
the certified floor drives Bob's LPIPS to $0.63$--$0.87$, leaving little usable
signal for the legitimate receiver; when the private subspace is task-null
(decoupled) the same floor costs almost nothing until
$\tau$ approaches the ceiling $\tau_{\max}(10^{4}){=}0.925$. The dotted line
is the certified floor itself, which the surrogate $P_e$ must dominate
(Theorem~\ref{th:certificate}).}
\label{fig:certificate}
\end{figure}

% ---------------------------------------------------------------------
% Fig. 4 -- leakage budget I_max(tau, K) sensitivity to the gallery size
% ---------------------------------------------------------------------
\begin{figure}[t]
\centering
\begin{tikzpicture}
  \pgfplotstableread{\figdatadir imax_gallery.dat}\ig
  \begin{axis}[
      width=\columnwidth, height=0.62\columnwidth,
      xlabel={certified error floor $\tau$},
      ylabel={$I_{\max}$ (bit)},
      xmin=0.30, xmax=0.92, ymin=0,
      tick label style={font=\footnotesize},
      label style={font=\footnotesize},
      legend style={font=\scriptsize, draw=none, fill=none,
                    at={(0.98,0.98)}, anchor=north east},
      ymajorgrids, grid style={densely dotted, black!25},
    ]
    \addplot[blue!55!black, dashed] table[x=tau, y=k1e3] {\ig};
    \addplot[blue!70!black, mark=*, mark size=1.2pt] table[x=tau, y=k1e4] {\ig};
    \addplot[orange!80!black, dotted, very thick] table[x=tau, y=k1e6] {\ig};
    \legend{$K=10^{3}$, $K=10^{4}$, $K=10^{6}$}
  \end{axis}
\end{tikzpicture}
\caption{Leakage budget $I_{\max}(\tau,K)$ of \eqref{eq:imax} versus the
certified floor, for three identity-gallery sizes. Raising $K$ pushes the
feasibility ceiling $\tau_{\max}(K)$ and the budget upward only
logarithmically: adjacent gallery sizes differ by about $0.025$ in
$\tau_{\max}$ and the full span over three orders of magnitude ($10^{3}$ to
$10^{6}$) is about $0.050$. Curves are truncated at their own feasibility
ceiling $\tau_{\max}(K)$; beyond it no non-negative budget exists.}
\label{fig:imax}
\end{figure}

% ---------------------------------------------------------------------
% Fig. 3 -- composite, one column per story
% ---------------------------------------------------------------------
\begin{figure*}[t]
\centering
\begin{tikzpicture}
  \pgfplotstableread{\figdatadir waterfill_sweep.dat}\wf
  \begin{axis}[
      width=0.315\textwidth, height=0.27\textwidth,
      xlabel={dual variable $\nu$}, ylabel={$\sigma_h^{2}/g_h$},
      xmode=log, ymode=log, ymin=1e-4, ymax=1e8,
      tick label style={font=\footnotesize},
      label style={font=\footnotesize},
      title={(a) water-filling}, title style={font=\footnotesize},
      legend style={font=\scriptsize, draw=none, fill=none,
                    at={(0.98,0.03)}, anchor=south east},
    ]
    \addplot[blue!55!black] table[x=mu, y=w_1] {\wf};
    \addplot[blue!55!black, densely dashed] table[x=mu, y=w_8] {\wf};
    \addplot[orange!80!black, densely dotted, very thick]
      table[x=mu, y=w_125] {\wf};
    \legend{$w_h$, $w_h/8$, $w_h/125$}
  \end{axis}
\end{tikzpicture}
\hfill
\begin{tikzpicture}
  \pgfplotstableread{\figdatadir accounting_audit.dat}\ac
  \begin{axis}[
      width=0.315\textwidth, height=0.27\textwidth,
      xlabel={noise-to-signal ratio (dB)}, ylabel={certified $\epsilon$},
      ymode=log,
      tick label style={font=\footnotesize},
      label style={font=\footnotesize},
      title={(b) audited $\epsilon$ at one release},
      title style={font=\footnotesize},
      legend style={font=\scriptsize, draw=none, fill=none,
                    at={(0.98,0.03)}, anchor=south east},
    ]
    \addplot[draw=none, fill=red!7] coordinates {
      (9.7,1)(24,1)(24,5000)(9.7,5000)(9.7,1)};
    \node[font=\scriptsize, red!60!black, anchor=north west]
      at (axis cs:10.2,3000) {hazard: $\mu_h\gg c$};
    \addplot[blue!55!black, mark=triangle*, mark size=1.5pt]
      table[x=nsr_db, y=eps_quadrature] {\ac};
    \addplot[green!40!black, mark=square*, mark size=1.2pt]
      table[x=nsr_db, y=eps_l1_exact] {\ac};
    \addplot[orange!80!black, mark=*, mark size=1.2pt, densely dashed]
      table[x=nsr_db, y=eps_l1_stylized] {\ac};
    \legend{quadrature, $\ell_1$ exact, $\ell_1$ ampl.}
  \end{axis}
\end{tikzpicture}
\hfill
\begin{tikzpicture}
  \pgfplotstableread{\figdatadir certificate_gap.dat}\cg
  \begin{axis}[
      width=0.315\textwidth, height=0.27\textwidth,
      xlabel={certified floor $\tau$}, ylabel={$P_e^{\rm sim}-\tau$},
      ymin=-0.05, ymax=0.62,
      tick label style={font=\footnotesize},
      label style={font=\footnotesize},
      title={(c) certificate tightness},
      title style={font=\footnotesize},
    ]
    \addplot[purple!60!black, mark=*, mark size=1.2pt]
      table[x=tau, y=gap] {\cg};
    \addplot[black, densely dotted, domain=0.42:0.93] {0};
  \end{axis}
\end{tikzpicture}
\caption{(a) Water-filling charge: one latent direction, three task weights.
Privacy is cheaper for a task-null direction because $a_h=w_h\kappa_h$
enters the KKT solution \eqref{eq:kkt}. (b) Audit of the accounting laws on
the \emph{same} release at the low-noise (large-$\mu$) operating point certified at
$\varepsilon{=}2000$, $\delta{=}10^{-5}$: quadrature (GDP) sits between the
rigorous tail bound and the amplitude-only $\ell_1$ form, the latter under-
reporting $\varepsilon$ once $\mu_h\gg c$. The red shaded band marks this
large-$\mu$ hazard region, where a Laplace-style $\ell_1$ composition yields an
assessment tighter than the guarantee. (c) The surrogate $P_e$ dominates the
certified floor everywhere: the certificate is conservative, not optimistic.}
\label{fig:composite}
\end{figure*}

% ---------------------------------------------------------------------
% Fig. 4 -- latent-level measurement on the released SemanticStyleGAN
% ---------------------------------------------------------------------
\begin{figure*}[t]
\centering
\begin{tikzpicture}
  \pgfplotstableread{\figdatadir latent_utility_landscape.dat}\lu
  \begin{axis}[
      width=0.315\textwidth, height=0.27\textwidth,
      ybar, bar width=3.2pt,
      ylabel={measured $w_h$},
      ymin=0, ymax=0.40,
      xtick={0,...,14},
      xticklabels={sh0,sh1,bg,skin,eye,brw,mth,nos,ear,hair,nck,clth,eyg,hat,erR},
      x tick label style={rotate=55, anchor=east, font=\tiny},
      tick label style={font=\footnotesize},
      label style={font=\footnotesize},
      title={(a) measured utility landscape}, title style={font=\footnotesize},
    ]
    \addplot[fill=blue!45, draw=blue!55!black] table[x=h, y=w_h] {\lu};
  \end{axis}
\end{tikzpicture}
\hfill
\begin{tikzpicture}
  \pgfplotstableread{\figdatadir latent_sweep_examples.dat}\ls
  \begin{axis}[
      width=0.315\textwidth, height=0.27\textwidth,
      xlabel={injected $\sigma_h^{2}$}, ylabel={Bob LPIPS},
      xmode=log,
      tick label style={font=\footnotesize},
      label style={font=\footnotesize},
      title={(b) per-block noise sweeps}, title style={font=\footnotesize},
      legend style={font=\scriptsize, draw=none, fill=none,
                    at={(0.03,0.97)}, anchor=north west},
    ]
    \addplot[blue!55!black, mark=*, mark size=1.2pt]
      table[x=sigma2, y=shared0] {\ls};
    \addplot[orange!80!black, mark=square*, mark size=1.2pt]
      table[x=sigma2, y=skin] {\ls};
    \addplot[green!40!black, mark=triangle*, mark size=1.5pt]
      table[x=sigma2, y=hair] {\ls};
    \addplot[purple!60!black, mark=diamond*, mark size=1.5pt, densely dashed]
      table[x=sigma2, y=eyes] {\ls};
    \legend{shared\_0, skin, hair, eyes}
  \end{axis}
\end{tikzpicture}
\hfill
\begin{tikzpicture}
  \pgfplotstableread{\figdatadir latent_kappa.dat}\lk
  \begin{axis}[
      width=0.315\textwidth, height=0.27\textwidth,
      xlabel={injected $\sigma_h^{2}$}, ylabel={residual fraction $\kappa_h$},
      xmode=log, ymin=0, ymax=1.05,
      tick label style={font=\footnotesize},
      label style={font=\footnotesize},
      title={(c) Wiener deprotector residual}, title style={font=\footnotesize},
      legend style={font=\scriptsize, draw=none, fill=none,
                    at={(0.97,0.97)}, anchor=north east},
    ]
    \addplot[blue!55!black, mark=*, mark size=1.2pt]
      table[x=sigma2, y=kappa_empirical] {\lk};
    \addplot[red!70!black, densely dashed, very thick]
      table[x=sigma2, y=kappa_analytic] {\lk};
    \legend{empirical (held-out), analytic}
  \end{axis}
\end{tikzpicture}
\caption{Latent-level measurement on the released SemanticStyleGAN generator.
(a) Measured per-block utility weights $w_h$: the perceptual cost of latent
noise concentrates on the shared codes, hair, skin and background, while the
facial-feature blocks (eyes, nose, ear\_r) are nearly free to protect. Block
order: shared\_0, shared\_1, background, skin, eyes, brows, mouth, nose, ears,
hair, neck, cloth, eye\_g, hat, ear\_r. (b) Example decode-and-LPIPS sweeps
with the quadratic fit $\mathrm{LPIPS}\approx\mathrm{lpips}_0+w_h\sigma_h^2$;
the eyes curve is almost flat---the generator absorbs local latent noise into
plausible variations. (c) Residual fraction of the per-block Wiener
deprotector on held-out codes versus the analytic prediction; the two agree to
three decimals, so $\kappa_h$ is a measured quantity, not a fit artefact.}
\label{fig:latent}
\end{figure*}


% =====================================================================
%  section_privacy_accounting.tex
%  Drop-in section for an IEEEtran two-column paper.
%  If your class already defines the theorem environments, delete the
%  three \newtheorem lines below (keep \theoremstyle).
% =====================================================================
\theoremstyle{plain}
\newtheorem{theorem}{Theorem}
\newtheorem{proposition}{Proposition}
\newtheorem{lemma}{Lemma}
\newtheorem{corollary}{Corollary}
\theoremstyle{definition}
\newtheorem{assumption}{Assumption}
\newtheorem{remark}{Remark}

\section{Privacy Accounting for Semantic Latents:
         Certificates, Allocation, and Invertible Shaping}
\label{sec:accounting}

\IEEEPARstart{T}{he} scheme studied in the companion
report~\cite{arXiv2510} perturbs a subset of the semantic latent codes and
drives the perturbation with a discriminator so that the released noise is
claimed to be statistically close to a Laplace draw while remaining
semantically mild. Two
parts of that construction are left implicit, and both affect how the claims
can be stated for publication: what exactly is being certified, and how the
certification cost enters the rate--distortion budget. This section supplies
the missing parts. We (i) isolate the component-wise sensitivity of a
decomposed semantic latent (Proposition~\ref{pr:inflation}), (ii) state
precisely what an invertible secret shaping can and cannot buy
(Theorem~\ref{th:postproc}), (iii) convert the informal ``distribution
resembles Laplace'' argument into a mutual-information certificate that yields
an operational lower bound on an adversary's identification error
(Theorem~\ref{th:certificate}), (iv) derive the allocation law that a given
accounting rule induces (Theorems~\ref{th:allocation}
and~\ref{th:shape}), and (v) give the complete
$\RDP\!\to\!(\varepsilon,\delta)$ and Gaussian-DP composition derivations
(Section~\ref{sec:rdpgdp}).

\subsection{Setup and Notation}
\label{sec:setup}

\emph{Semantic latent.} A semantic encoder (here a StyleGAN generator inverted
by a Semantic StyleGAN encoder~\cite{semanticstylegan}) maps an image to a
latent $z\in\mathbb{R}^{n}$ decomposed into $H$ semantic components
\begin{equation}
 z=(z_1,\dots,z_H),\qquad z_h\in\mathbb{R}^{d_h},\qquad
 n=\textstyle\sum_{h=1}^{H} d_h ,
\label{eq:decomp}
\end{equation}
where $h$ indexes a shared (global) code or a local (part-wise) code group.
For the reference backbone, $H_{\mathrm c}=28$ codes of dimension
$d_h=512$, $n=14336$; our implementation groups the 13 part-wise local codes
into 13 blocks of $d_h=1024$ plus 2 shared blocks of $d_h=512$ ($H=15$,
$n=14336$). Grouping affects only bookkeeping; the composition formulas of
Section~\ref{sec:rdpgdp} are additive per component and are unchanged by it.

\emph{Private subspace.} Let $P\subseteq\{1,\dots,H\}$ be the set of
\emph{declared private} components and $Q=[H]\setminus P$ its complement.
$P$ is a design decision, not a derived quantity: it is the set of semantic
attributes whose inference the system promises to obstruct (identity, in the
reference task). We write $z_P=(z_h)_{h\in P}$.

\begin{assumption}[Bounded coordinate range]
\label{as:range}
Every latent coordinate satisfies $z_{h,i}\in[a_h,b_h]$; without loss of
generality $b_h-a_h=b-a$ for all $(h,i)$, after the $1$-Lipschitz
normalisation applied by the encoder.
\end{assumption}

\begin{assumption}[Block-wise second-order structure]
\label{as:indep}
For each block $h$, the coordinate variances are uniformly bounded,
$\operatorname{Var}(z_{h,i})\le g_{h,i}$ for every coordinate $i$, and the
block source power satisfies
$\operatorname{Tr}\Sigma_h=\sum_i\operatorname{Var}(z_{h,i})\le g_h$. No PSD
ordering $\Sigma_h\preceq\operatorname{diag}(g_{h,i})$ is assumed; residual
within-block and across-block correlation is charged conservatively to the
block source power $g_h$, see Remark~\ref{rm:corr}.
\end{assumption}

\emph{Vector-release convention.} Each block $h$ is released as a
$d_h$-dimensional vector. The Gaussian worst-case mutual information uses the
relaxed trace bound
$I(z_h;z_h+N_h)\le\tfrac12\log(1+\operatorname{Tr}\Sigma_h/\sigma_h^2)$, so
$g_h$ denotes the block source power $\operatorname{Tr}\Sigma_h$; the
coordinate count $d_h$ enters both the trace proxy and the block $\ell_2$
sensitivity $\Delta f_h=\sqrt{d_h}(b-a)$.

\emph{Adjacency and code-level change.} Two records are adjacent, $D\sim D'$,
if they differ in the private semantics of exactly one subject (replace-one).
A local block concatenates two 512 codes; if adjacency changes only one of
them, the true $\ell_2$ sensitivity is $(b-a)\sqrt{512}$. The sensitivity
used below takes the conservative worst case $\Delta f_h=(b-a)\sqrt{d_h}$,
over-estimating a partially changed block by at most $\sqrt{2}$. All
quantities are per-record.

\emph{Release mechanism.} The transmitter releases
\begin{equation}
Y=T\!\left(z+N\right),\qquad
N_h\sim\mathcal{N}\!\left(0,\sigma_h^{2} I_{d_h}\right)\ \text{independent},
\label{eq:mechanism}
\end{equation}
where $T$ is an optional transformation. The baseline scheme
of~\cite{arXiv2510} instantiates $T$ as a learned, per-component shaping
network trained against a discriminator, with the additional requirement that
$T$ be invertible at the legitimate receiver. The legitimate receiver observes
$Y$; the eavesdropper observes $Y$. In the worst case $T$ is invertible and
known to the eavesdropper, so $Y=Y$, which is the case certified below.

\subsection{Sensitivity Accounting}
\label{sec:sensitivity}

\begin{proposition}[Accounting, not mechanism, inflates the noise]
\label{pr:inflation}
Under Assumption~\ref{as:range}, the L2 sensitivity of the whole latent and of
a single component are
\begin{equation}
\Delta f_{\mathrm{glob}}=(b-a)\sqrt{{\textstyle\sum_h d_h}},
\qquad
\Delta f_h=(b-a)\sqrt{d_h}.
\label{eq:sens}
\end{equation}
Three bookkeeping rules follow:
\begin{enumerate}
\item[(i)] \emph{per-component with $\ell_1$ composition:}
$\varepsilon=\sum_{h}\Delta f_h/\sigma_h$, so the smallest equalising noise
satisfies $\sigma\ge(b-a)\textstyle\sum_h\sqrt{d_h}/\varepsilon$;
\item[(ii)] \emph{monolithic vector query:}
$\sigma\ge\Delta f_{\mathrm{glob}}/\varepsilon
=(b-a)\sqrt{{\textstyle\sum_h d_h}}/\varepsilon$;
\item[(iii)] \emph{charging every component the global sensitivity} (the
practice implied when one reported $\varepsilon$ is reused for every code):
$\sigma\ge H_{\mathrm c}\,\Delta f_{\mathrm{glob}}/\varepsilon$.
\end{enumerate}
For $d_h\equiv d$, the ratios are
$\sigma_{\mathrm{(i)}}/\sigma_{\mathrm{(ii)}}=\sqrt{H_{\mathrm c}}$ and
$\sigma_{\mathrm{(iii)}}/\sigma_{\mathrm{(i)}}=\sqrt{H_{\mathrm c}}$; the
third-to-first factor is $H_{\mathrm c}$ in noise power, or
$10\log_{10}H_{\mathrm c}$ decibel.
\end{proposition}

\begin{proof}
(i) is the standard Gaussian-mechanism bound applied blockwise, using
$\|\Delta z_h\|_2\le(b-a)\sqrt{d_h}$ on every adjacent pair.
(ii) follows from $\|\Delta z\|_2^2=\sum_h\|\Delta z_h\|_2^2
\le(b-a)^2\sum_h d_h$. (iii) substitutes $\Delta f_{\mathrm{glob}}$ for every
summand in (i). For $d_h\equiv d$:
$\sum_h\sqrt d=H_{\mathrm c}\sqrt d$ and
$\sqrt{\sum_h d}=\sqrt{H_{\mathrm c}d}$, giving
$\sigma_{\mathrm{(i)}}/\sigma_{\mathrm{(ii)}}=\sqrt{H_{\mathrm c}}$, and
$\sigma_{\mathrm{(iii)}}/\sigma_{\mathrm{(i)}}=
H_{\mathrm c}\sqrt{H_{\mathrm c}d}/(H_{\mathrm c}\sqrt d)=\sqrt{H_{\mathrm c}}$.
\end{proof}

With $H_{\mathrm c}=28$ the third-to-first factor is $\sqrt{28}=5.29$ in
amplitude and $28\times$ ($14.5$\,dB) in power. This is the full extent of
the available gain, and it should be described accurately: it is a correction
of bookkeeping, not a new mechanism. A correct monolithic vector accounting
(ii) is already as tight as correct per-component quadrature accounting
(Proposition~\ref{pr:quadrature}). Any larger dividend must come from the
allocation itself (Section~\ref{sec:allocation}), which is a separate degree
of freedom.

\subsection{What Shaping Can and Cannot Do}
\label{sec:shaping}

\begin{theorem}[Post-processing and invertible shaping]
\label{th:postproc}
Fix the noise law in \eqref{eq:mechanism}.
\begin{enumerate}
\item[(i)] For every $T$ that is a (possibly randomised) function of the
mechanism output with randomness independent of $D$, $T\circ Y$ is
$(\varepsilon,\delta)$-DP whenever $Y$ is.
\item[(ii)] If $T$ is invertible and $T^{-1}$ is available to the adversary,
$T\circ Y$ provides exactly the same guarantee as $Y$: the adversary recovers
$Y$ and the protection, if any, is due to $N$ alone.
\item[(iii)] If $T$ depends on the private latent directly rather than on the
release, it is part of the mechanism and its sensitivity must be re-accounted;
$T$ may then \emph{increase} $\varepsilon$.
\end{enumerate}
\end{theorem}

\begin{proof}
(i) is the post-processing inequality: for adjacent $D\sim D'$ there is a
coupling $(Y,Y')$ of the two output laws with $\Pr[Y\neq Y']\le\delta$;
applying the same $T$ to both sides preserves
$\{T(Y)=T(Y')\}\supseteq\{Y=Y'\}$, hence
$\Pr[T(Y)\neq T(Y')]\le\delta$, and on the complement the two conditional
laws coincide, which is exactly $(\varepsilon,\delta)$-DP.
(ii) Eve computes $T^{-1}$, so the information available to her is the
$\sigma$-field generated by $Y$; by the data-processing inequality
$I(z;T^{-1}(T(Y)))=I(z;Y)$ and every distinguishing advantage is unchanged.
(iii) $T(z+\cdot)$ is a function of the data, so the pair
$(z,T(z+N))$ is released by a mechanism whose sensitivity is governed by
$\|T(z)-T(z')\|$, not by $\Delta f$. A per-image contrast normalisation
$\sigma\mapsto\sigma/\widehat{s}(z)$ therefore leaks through
$\widehat{s}(z)$.
\end{proof}

\begin{remark}[The invertibility--privacy tension]
\label{rm:tension}
Part~(ii) imposes a constraint that applies to any deployment: DP noise cannot
be
``learned away'' for the legitimate receiver only. The deprotection network at
Bob is a decoder of the released signal; if it recovers $z$ up to a residual
that is useful to Bob, the same architecture applied by Eve recovers the same
residual. Hence Bob's advantage over Eve cannot originate in denoising a
differentially private release. It can only originate in
(a) a \emph{secret} shaping $T$ that Eve does not know, in which case the
guarantee is key-conditional and must be stated as such; it is not
$(\varepsilon,\delta)$-DP in the Kerckhoffs sense; or
(b) noise whose projection onto the task subspace is zero, in which case Bob
needs no denoiser and the protocol reduces to an allocation over the
components indexed by $P$ that actually project onto the task subspace.
The baseline scheme of~\cite{arXiv2510} relies on a learned noise pattern to
keep Bob's distortion low at large $\varepsilon$; by
Theorem~\ref{th:postproc} and this remark that mechanism can only be justified
as (a), i.e.\ as a keyed component. A claim of pure DP protection together
with negligible legitimate distortion at large $\varepsilon$ would then be
mis-attributed.
\end{remark}

\subsection{A Certificate for the Private Subspace}
\label{sec:certificate}

Data processing bounds the leakage of any decision Eve runs on the release.
We use this to convert a privacy budget into an operational error bound for the
attack family the system cares about: identification among $K$ enrolled
templates.

\begin{lemma}[Gaussian worst case at fixed covariance]
\label{lm:maxent}
Let $z_h\in\mathbb{R}^{d_h}$ be the released block with covariance $\Sigma_h$
and block source power $g_h=\operatorname{Tr}\Sigma_h$, and let
$N_h\sim\mathcal{N}(0,\sigma_h^{2}I_{d_h})$ be independent of $z_h$. Then
\begin{equation}
I(z_h;z_h+N_h)\le\tfrac12\log\det\!\left(I+\frac{\Sigma_h}{\sigma_h^{2}}\right)
\le\tfrac12\log\!\left(1+\frac{g_h}{\sigma_h^{2}}\right).
\label{eq:maxent}
\end{equation}
\end{lemma}

\begin{proof}
Write $Y_h=z_h+N_h$. Among random vectors with fixed covariance $\Sigma_h$, the
Gaussian achieves the maximal differential entropy, so
$h(z_h)\le\tfrac12\log\big((2\pi e)^{d_h}\det\Sigma_h\big)$ and
$h(Y_h)\le\tfrac12\log\big((2\pi e)^{d_h}\det(\Sigma_h+\sigma_h^{2}I)\big)$;
the conditional entropy is
$h(Y_h\mid z_h)=h(N_h)=\tfrac12\log\big((2\pi e\sigma_h^{2})^{d_h}\big)$. Hence
\begin{equation*}
I(z_h;Y_h)=h(Y_h)-h(Y_h\mid z_h)
\le \tfrac12\log\det\!\left(I+\frac{\Sigma_h}{\sigma_h^{2}}\right).
\end{equation*}
For a PSD matrix $A$, $\det(I+A)\le 1+\operatorname{Tr}A$, giving the relaxed
trace bound $\tfrac12\log(1+\operatorname{Tr}\Sigma_h/\sigma_h^{2})$, which is
\eqref{eq:maxent}.
\end{proof}

\begin{theorem}[Leakage certificate and identification floor]
\label{th:certificate}
Grant Assumptions~\ref{as:range}--\ref{as:indep}, let the private template
$B=f(z_P)$ take $K$ values, and let the adversary observe
$Y=T(z+N)$ with $T$ invertible (the worst case). If
\begin{equation}
I_{\mathrm{leak}}(\sigma)
=\tfrac12\!\!\sum_{h\in P}\!\log\!\left(1+\frac{g_h}{\sigma_h^{2}}\right)
\le I_{\max}(\tau),
\label{eq:leak}
\end{equation}
where
\begin{equation}
I_{\max}(\tau)=\log_2 K-1-\tau\log_2(K-1)\quad[\text{bit}],
\label{eq:imax}
\end{equation}
then every estimator $\widehat B(Y)$ has identification error
$\Pr[\widehat B\neq B]\ge\tau$.
\end{theorem}

\begin{proof}
The Markov chain $B\to z_P\to Y=T(z+N)\to z+N$ (the last step is the
deterministic $T^{-1}$) gives, by data processing,
$I(B;Y)\le I(z_P;z+N)$. Under the block-wise second-order
structure of Assumption~\ref{as:indep} and Lemma~\ref{lm:maxent},
\begin{equation}
I(z_P;z+N)\le\tfrac12\!\!\sum_{h\in P}\!\log\!\left(1+\frac{g_h}{\sigma_h^{2}}\right)
= I_{\mathrm{leak}}(\sigma),
\label{eq:cert-bound}
\end{equation}
the right-hand side of \eqref{eq:leak}. Fano's inequality for a $K$-ary $B$~\cite{fano}
gives
\begin{equation}
\Pr[\widehat B\neq B]\ge\frac{H(B\mid Y)-1}{\log_2(K-1)}
\ge\frac{\log_2 K-I(B;Y)-1}{\log_2(K-1)},
\end{equation}
so $I(B;Y)\le I_{\max}(\tau)$ implies the claimed floor. Sending
$I_{\mathrm{leak}}\downarrow0$ in \eqref{eq:imax} yields the feasibility
condition
\begin{equation}
\tau<\tau_{\max}(K)=\frac{\log_2K-1}{\log_2(K-1)}.
\label{eq:taumax}
\end{equation}
$\tau\ge\tau_{\max}$ would require a negative information budget, which is
infeasible.
\end{proof}

The Fano bound above uses a uniform prior on $B$. For a non-uniform prior the
entropy term is $H(B)\le\log_2K$, so the right-hand side weakens by
$\log_2K-H(B)$; the uniform case is the strongest guarantee available, and a
measured prior should be inserted instead of $\log_2K$ in a deployed system.

\begin{corollary}[Numbers]
\label{co:numbers}
$\tau_{\max}(10^{4})=0.9248$, $\tau_{\max}(10^{6})=0.9498$, and
$\tau_{\max}(K)\uparrow1$ only logarithmically: an arbitrarily strong
\emph{certified} identification guarantee is unattainable for any finite
template bank, because $H(B)=\log_2K$ is itself the currency. In our
calibration, a floor $\tau=0.90$ with $K=10^{4}$ leaves
$I_{\max}=0.33$\,bit, the entire budget for all private blocks jointly.
\end{corollary}

\begin{remark}[Scope of the certificate]
\label{rm:scope}
Theorem~\ref{th:certificate} certifies the identification family, not
arbitrary inference: it constrains every estimator of $B$, but says nothing
about an attribute $A$ that is not a function of $z_P$. It is a
task-level guarantee, strictly weaker than DP and strictly stronger than the
empirically measured reduction in an attack's identification accuracy.
Both must be reported, and neither may be presented as the other.
\end{remark}

\begin{remark}[Residual correlation across blocks]
\label{rm:corr}
Hadamard's inequality for a PSD block matrix gives
$\det(I+\Sigma_{\mathrm{joint}})\le\prod_h\det(I+\Sigma_{hh})$. Hence the
block-wise marginal bound in \eqref{eq:cert-bound} remains an upper bound on
the joint leakage even when blocks are correlated: off-diagonal entries of
$\Sigma_{\mathrm{joint}}$ only make the Hadamard gap larger, which loosens
the bound but does not invalidate it. A joint covariance estimate would be
needed to tighten the bound; charging residual correlation to the block
variances $g_h$ is the conservative way to retain soundness without that
estimate.
\end{remark}

\subsection{The Allocation Law Induced by the Accounting Rule}
\label{sec:allocation}

Write $x_h=\sigma_h^{2}$ and $a_h=w_h\kappa_h>0$, where $w_h$ is the
calibrated sensitivity of the task loss to the block and $\kappa_h$ maps
latent power to distortion. The transmitter solves
\begin{equation}\begin{split}
\min_{x_h\ge0}\ &\sum_{h=1}^{H} a_h x_h\\
\text{s.t.}\ &\sum_{h\in P}\varphi_h(x_h)\le B,\quad
\varphi_h(x)=\tfrac12\log\!\left(1+\frac{g_h}{x}\right),
\end{split}\label{eq:primal}
\end{equation}
with $B$ the allowed leakage in nats. Noise is placed only on $P$: a block in
$Q$ carries no obligation and no benefit, so its optimum is $x_h=0$.

\begin{theorem}[Convexity, uniqueness, and closed form]
\label{th:allocation}
The problem \eqref{eq:primal} is a convex program with a unique optimum when
$B\in(0,\infty)$. Let $\nu\ge0$ denote its dual multiplier. The KKT solution is
\begin{equation}
x_h^{\star}(\nu)=\frac{g_h}{2}\left(
\sqrt{1+\frac{2\nu}{a_h g_h}}-1\right),\qquad h\in P,
\label{eq:kkt}
\end{equation}
pinned by $\sum_{h\in P}\varphi_h(x_h^\star(\nu))=B$. The map
$\nu\mapsto\sum_{h\in P}\varphi_h(x_h^\star(\nu))$ is strictly decreasing and
continuous from $+\infty$ ($\nu\downarrow0$) to $0$ ($\nu\uparrow\infty$), so
the bisection in the accompanying code is well posed and converges
geometrically in $\log\nu$.
\end{theorem}

\begin{proof}
$\varphi_h$ is convex on $x>0$:
$\varphi_h'(x)=-\frac{g_h}{2x(x+g_h)}$ and
$\varphi_h''(x)=\frac{g_h(2x+g_h)}{2x^2(x+g_h)^2}>0$; the objective is linear,
so the feasible set is convex and the program is convex. At an optimum the
constraint is active because $B>0$ forces some $x_h<\infty$, and on the active
set every $x_h>0$ for $h\in P$; strict convexity of $\varphi_h$ then gives a
unique minimiser. Stationarity reads
$a_h+\nu\,\varphi_h'(x_h)=0$, i.e.
$a_h=\nu g_h/(2x_h(x_h+g_h))$, which rearranges to
$x_h^2+g_hx_h-\nu g_h/(2a_h)=0$; its positive root is \eqref{eq:kkt}.
Monotonicity: $\partial x_h^\star/\partial\nu>0$ for every $h$, and
$\varphi_h$ is decreasing in $x$, so each summand decreases strictly; the
limits follow from $x_h^\star\downarrow0$ and $x_h^\star\uparrow\infty$. A
continuous strictly monotone scalar function is inverted by bisection, and
each halving of the $\log\nu$ interval halves the bracketing width, whence
geometric convergence.
\end{proof}

\begin{corollary}[Two asymptotics]
\label{co:asymptotics}
From \eqref{eq:kkt},
\begin{align}
x_h^\star&\simeq\frac{\nu}{2a_h}, & x_h&\ll g_h,\label{eq:asym1}\\
x_h^\star&\simeq\sqrt{\frac{\nu g_h}{2a_h}}, & x_h&\gg g_h.\label{eq:asym2}
\end{align}
The first regime is task-driven and independent of $g_h$; the second is
geometry-driven and grows only as $\sqrt{g_h}$. A component that is cheap for
the task but expensive to protect ($a_h$ small, $g_h$ large) absorbs most of
the noise precisely in the second regime.
\end{corollary}

\begin{theorem}[The accounting rule dictates the shape]
\label{th:shape}
Replace the leakage constraint of \eqref{eq:primal} by each of the standard
accounting laws at a common budget $\varepsilon$, and define
\begin{equation}
c=\sqrt{2\log(1.25/\delta)}.
\label{eq:tailconst}
\end{equation}
The stationary profiles are:
\begin{enumerate}
\item[(a)] \emph{amplitude-only $\ell_1$ composition}
$\sum_h c\Delta f_h/\sigma_h\le\varepsilon$:
\begin{equation}
x_h^{\star}\propto\Big(\frac{\Delta f_h}{a_h}\Big)^{2/3},
\label{eq:l1shape}
\end{equation}
a geometric profile, strongly unequal and dominated by components with large
sensitivity and small task weight.
\item[(b)] \emph{quadrature composition} (RDP and GDP alike,
Section~\ref{sec:rdpgdp}): the profile \eqref{eq:kkt} with asymptotics
\eqref{eq:asym1}--\eqref{eq:asym2},
a smooth logarithmic profile whose ratios depend only on $a_h/a_{h'}$ in the
task-driven regime and additionally on $g_h/g_{h'}$ in the geometry-driven
regime.
\item[(c)] \emph{rigorous tail bound} (Equation~\eqref{eq:l1exact} below):
\begin{equation}
2a_hx_h^{2}=\nu\big(\Delta f_h^{2}+c\Delta f_h\sqrt{x_h}\big),
\label{eq:l1exactshape}
\end{equation}
so $x_h\propto\Delta f_h/\sqrt{a_h}$ when $\mu_h\gg c$ (the quadratic term
dominates) and \eqref{eq:l1shape} when $\mu_h\ll c$ (the linear term dominates),
with $\mu_h=\Delta f_h/\sigma_h$.
\end{enumerate}
\end{theorem}

\begin{proof}
(a) The constraint is $\sum_hb_hx_h^{-1/2}\le\varepsilon$ with
$b_h=c\Delta f_h$; it is convex since $\psi_h(x)=b_hx^{-1/2}$ has
$\psi_h''>0$. Stationarity $a_h=\nu b_hx_h^{-3/2}/2$ gives
$x_h=(2a_h/(\nu b_h))^{-2/3}\propto(b_h/a_h)^{2/3}$, which is
\eqref{eq:l1shape}; $\nu$ is fixed by the budget.
(b) GDP accumulates in quadrature,
$\sum_h(\Delta f_h/\sigma_h)^{2}\le c_{\mathrm{GDP}}(\varepsilon)$ by
\eqref{eq:gdpsum} inverted in $\varepsilon$; RDP accumulates
$\sum_h\alpha\Delta f_h^2/(2\sigma_h^2)\le\varepsilon$ at the selected order
$\alpha$ by \eqref{eq:rdpsum}. Both are quadratic in $1/\sigma_h$, hence of
the form $\sum_h\chi_h/x_h\le B'$ after absorbing $\Delta f_h^2$, which is the
same KKT system as \eqref{eq:primal} with a rescaled multiplier; the stated
dependence of $x_h/x_{h'}$ follows from \eqref{eq:kkt}.
(c) The constraint \eqref{eq:l1exact} has
$f_h(x)=\Delta f_h^2/(2x)+c\Delta f_hx^{-1/2}$; stationarity
$a_h=\nu f_h'(x_h)$ reads
$2a_hx_h^2=\nu\Delta f_h(\Delta f_h+c\sqrt{x_h})$, which is
\eqref{eq:l1exactshape}. Dividing by $c\Delta f_h\sqrt{x_h}$ identifies the
governing ratio $\mu_h/c=\Delta f_h/(c\sqrt{x_h})$: for $\mu_h\gg c$ the term
$\Delta f_h^2$ dominates and $x_h^2\approx\nu\Delta f_h^2/(2a_h)$, giving
$x_h\propto\Delta f_h/\sqrt{a_h}$; for $\mu_h\ll c$ the term
$c\Delta f_h\sqrt{x_h}$ dominates and \eqref{eq:l1shape} is recovered.
\end{proof}

\begin{remark}[One exponent, three laws, one message]
\label{rm:exponents}
The three laws differ in the exponent of the task weight: $1$ for
\eqref{eq:kkt} in the task-driven regime, $2/3$ for the amplitude-only
$\ell_1$ law, $1/2$ for the rigorous tail bound at large $\mu$. All three agree
that noise belongs where $a_h$ is small; the disagreement is in how strongly
they push. The static-index policy has no exponent at all, because it cannot
use the $a_h$ scale.
\end{remark}

\begin{corollary}[Static index selection is a corner of this program]
\label{co:static}
The reference strategy --- protect a fixed set $S$ of $k$ codes with a common
noise level and leave the rest untouched --- is the feasible point of the
$\ell_1$ program obtained by imposing the extra constraints
$x_h=\sigma_{\max}^{2}\mathbf 1\{h\in S\}$,
$x_h=\sigma_{\min}^{2}\mathbf 1\{h\notin S\}$ with
$\sigma_{\min}\approx0$, $|S|=k$. These constraints define a subset of the
feasible set of the $\ell_1$ program in Theorem~\ref{th:shape}(a), so the
static rule cannot outperform
\eqref{eq:l1shape}, and it is strictly worse whenever the $\ell_1$ optimum is
supported on more than $k$ components with unequal amplitudes --- the generic
case, since $\Delta f_h$ and $a_h$ are heterogeneous across codes. The dynamic
allocation contains the static index rule as a degenerate special case and
replaces the combinatorial choice of $S$ by a global multiplier.
\end{corollary}

\subsection{RDP, Gaussian-DP, and Their Composition}
\label{sec:rdpgdp}

\emph{Gaussian mechanism, one block.} Let $q(z)=z$ have sensitivity
$\Delta_h=\Delta f_h$ on the block. For neighbouring $z,z'$ with
$\delta_h=z'-z$, $\|\delta_h\|_2\le\Delta_h$, the laws of $Y_h$ are
$\mathcal{N}(z_h,\sigma_h^2I)$ and $\mathcal{N}(z_h',\sigma_h^2I)$. The
Rényi divergence of order $\alpha>1$ between two such Gaussians is
\begin{equation}
\begin{aligned}
D_\alpha\!\left(\mathcal{N}(z_h,\sigma_h^2I)\,\|\,
\mathcal{N}(z_h',\sigma_h^2 I)\right)
&=\frac{\alpha\|\delta_h\|^2}{2\sigma_h^2}
\le\frac{\alpha\Delta f_h^{2}}{2\sigma_h^{2}}\\
&=:\varepsilon_h^{\RDP}(\alpha),
\end{aligned}
\label{eq:rdpblock}
\end{equation}
obtained by evaluating the $\alpha$-th moment of the density ratio under the
shifted law.

\emph{Composition (RDP).} Rényi divergence adds over independent mechanism
outputs~\cite{mironov2017rdp}; the block releases are independent by
\eqref{eq:mechanism}, so
\begin{equation}
\varepsilon^{\RDP}(\alpha)=\sum_{h\in P}\varepsilon_h^{\RDP}(\alpha)
=\frac{\alpha}{2}\sum_{h\in P}\frac{\Delta f_h^{2}}{\sigma_h^{2}}.
\label{eq:rdpsum}
\end{equation}
Equation~\eqref{eq:rdpsum} is the quadrature, not the $\ell_1$, accumulation,
and grouping of codes does not change it.

\emph{Conversion to $(\varepsilon,\delta)$.} An
$(\alpha,\varepsilon^{\RDP})$-RDP mechanism is
$(\varepsilon^{\RDP}+\log(1/\delta)/(\alpha-1),\delta)$-DP for every
$\alpha>1$~\cite{mironov2017rdp}; optimising over $\alpha$,
\begin{equation}
\varepsilon(\delta)=\min_{\alpha>1}\left[
\frac{\alpha}{2}\sum_{h\in P}\frac{\Delta f_h^{2}}{\sigma_h^{2}}
+\frac{\log(1/\delta)}{\alpha-1}\right].
\label{eq:rdpconv}
\end{equation}
The bracket is convex in $\alpha$, so the minimiser is unique; the
implementation in the accompanying code searches over an order grid. The
conversion is conservative: tighter conversions are available in the
hypothesis-testing privacy literature, and \eqref{eq:rdpconv} is kept so that
all reported $\varepsilon$ are upper bounds.

\emph{Privacy loss random variable and Gaussian-DP.} For one block the privacy
loss of an outcome $y$ is
\begin{equation}
L_h(y)=\log\frac{p_{z_h}(y)}{p_{z_h'}(y)}
=\frac{\|\delta_h\|^{2}-2\langle y-z_h,\delta_h\rangle}{2\sigma_h^{2}},
\label{eq:loss}
\end{equation}
using $p_{z'}(y)=p_z(y-\delta_h)$. Under $y\sim p_{z_h}$,
$u=y-z_h$ has mean zero and covariance $\sigma_h^2I$, so
\begin{equation}
L_h\sim\mathcal{N}\!\left(\frac{\mu_h^{2}}{2},\,\mu_h^{2}\right),
\qquad \mu_h:=\frac{\Delta f_h}{\sigma_h}.
\label{eq:gdpblock}
\end{equation}
The summands are independent, so two Gaussians compose into a Gaussian:
\begin{equation}\begin{split}
L&=\sum_{h\in P}L_h
\sim\mathcal{N}\!\left(\frac{\mu_{\mathrm{tot}}^{2}}{2},
\mu_{\mathrm{tot}}^{2}\right),\\
\mu_{\mathrm{tot}}&=\sqrt{\sum_{h\in P}\mu_h^{2}}
=\sqrt{\sum_{h\in P}\frac{\Delta f_h^{2}}{\sigma_h^{2}}}.
\end{split}\label{eq:gdpsum}
\end{equation}
This is Gaussian differential privacy~\cite{dong2022gdp}. The conservative
one-sided form used in the code and in all tables is
\begin{equation}
\delta(\varepsilon)\le\Phi\!\left(-\frac{\varepsilon}{\mu_{\mathrm{tot}}}
+\frac{\mu_{\mathrm{tot}}}{2}\right),
\label{eq:gdpdelta}
\end{equation}
whose inverse is
$\varepsilon(\delta)=\mu_{\mathrm{tot}}(\mu_{\mathrm{tot}}/2-\Phi^{-1}(\delta))$.
The exact two-sided expression subtracts the second tail term; it differs
from \eqref{eq:gdpdelta} only for $\varepsilon$ of order
$\mu_{\mathrm{tot}}^2/2$, where both forms are already vacuous, so
\eqref{eq:gdpdelta} is the standard choice for tables.

\emph{Linear composition (basic accounting).} The classical Gaussian-mechanism
tail bound, composed linearly~\cite{dwork2014}, keeps the quadratic privacy
term:
\begin{equation}
\varepsilon_{\ell_1}(\delta)=\sum_{h\in P}\Big(\mu_h^{2}/2+c\,\mu_h\Big),
\qquad \mu_h=\frac{\Delta f_h}{\sigma_h},
\label{eq:l1exact}
\end{equation}
with $c$ of \eqref{eq:tailconst}. The amplitude-only truncation
$\sum_hc\Delta f_h/\sigma_h$ drops the $\mu_h^{2}/2$ term; it is the law that
Laplace-style accounting effectively assumes. The dropped term is not cosmetic.
At the operating point of the reference scheme~\cite{arXiv2510},
$\sigma_h\approx3.6$,
$\Delta f_h\approx91$, hence $\mu_h\approx24$ and $|P|=6$, and $\mu_h\gg c=4.84$;
\eqref{eq:l1exact} is then dominated by $\mu_h^2/2$. Concretely, at the
release our quadrature accounting certifies as $\varepsilon=2000$, the
amplitude-only truncation reports $\varepsilon=693$ --- an unsound claim,
smaller than that of the very scheme it is compared with --- while
\eqref{eq:l1exact} reports $2441$.
Figure~\ref{fig:composite}(b) plots this audit over the whole range. This is
the failure mode of a ``the noise resembles Laplace, hence
$\varepsilon=\Delta f/s$'' argument: the ratio form is dimensionally plausible
and numerically unsound.

\begin{proposition}[Quadrature is no worse, and is the correct one]
\label{pr:quadrature}
For any $\{\sigma_h\}$ and any $P$,
$\sum_{h\in P}\Delta f_h^{2}/\sigma_h^{2}
\le\big(\sum_{h\in P}\Delta f_h/\sigma_h\big)^{2}$, with equality iff at most
one summand is non-zero. Hence for a fixed noise vector, quadrature accounting
certifies a smaller $\varepsilon$ than the amplitude-only linear bound at the
same $\delta$, and it is the accumulation law that matches the actual
additive-noise mechanism of \eqref{eq:mechanism}.
\end{proposition}

\begin{proof}
The inequality is $\|\cdot\|_2^2\le\|\cdot\|_1^2$ on the vector
$(\Delta f_h/\sigma_h)_{h\in P}$, with equality iff the vector is supported on
a single coordinate. The last claim is \eqref{eq:rdpsum}--\eqref{eq:gdpsum}
computed from the likelihood ratio of the actual release.
\end{proof}

\subsection{Complexity and What This Section Does Not Claim}

\emph{Complexity.} The closed form \eqref{eq:kkt} is $O(H)$; the multiplier
search is bisection over $\log\nu\in[-18,18]$ with 50 steps, i.e.\ $O(50H)$;
the $(\varepsilon,\delta)$ conversion costs $O(|\mathcal A|H)$ for
$|\mathcal A|=140$ orders in \eqref{eq:rdpconv}. Calibration is a single
backward-free pass over latents; the budget solve is dominated by neither
encoder nor inverter forward passes, and the exact FLOP ratio should be
reported against the measured backbone in a numerical study.

\emph{Non-claims.} For the record: (i) Theorem~\ref{th:certificate} is a
guarantee for identification, not for arbitrary inference; (ii) it assumes the
block-wise second-order structure of Assumption~\ref{as:indep}, and residual
correlation only loosens it (Remark~\ref{rm:corr}); (iii) the private
subspace $P$ is declared before the noise is drawn; (iv) the distortion model
$a_h\sigma_h^{2}$ is a Jacobian surrogate, not a measured LPIPS, and must be
replaced by the calibrated curve before any table is quoted; (v)
$\tau_{\max}(K)$ in \eqref{eq:taumax} is a real ceiling and should be
reported next to every certified floor; (vi) the shaping $T$ of
\eqref{eq:mechanism} does not raise any guarantee (Remark~\ref{rm:tension}) and
is legitimate only as a keyed component.


\subsection*{Code and Data Availability}
The calibration fixture, the closed-form allocator and its accounting
routines, the reproduction pipeline, the PGFPlots data files and the unit
tests are available at the public repository recorded in the submission
metadata (DOI to be inserted upon publication). All reported numbers are
reproduced by the included one-command script; no measured latent corpus is
shipped in this preview, and the numerical results are based on the synthetic
anchor calibration described in Section~\ref{sec:setup}.

\begin{thebibliography}{9}
\bibitem{dong2022gdp} J.~Dong, A.~Roth, and W.~J. Su,
``Gaussian differential privacy,'' \emph{J. Roy. Stat. Soc. B}, vol.~84,
no.~1, pp.~3--37, 2022.
\bibitem{mironov2017rdp} I.~Mironov, ``R\'enyi differential privacy,'' in
\emph{Proc. IEEE 30th Comput. Secur. Found. Symp. (CSF)}, 2017, pp.~263--275.
\bibitem{dwork2014} C.~Dwork and A.~Roth, \emph{The Algorithmic Foundations of
Differential Privacy}. Hanover, MA, USA: Now Publishers, 2014.
\bibitem{fano} T.~M. Cover and J.~A. Thomas, \emph{Elements of Information
Theory}, 2nd ed. Hoboken, NJ, USA: Wiley, 2006.
\bibitem{semanticstylegan} Y.~Shi \emph{et al.}, ``Semantic StyleGAN: Learning
compositional generative priors for controllable image synthesis and editing,''
in \emph{Proc. IEEE/CVF Conf. Comput. Vis. Pattern Recognit. (CVPR)}, 2022,
pp.~11254--11264.
\bibitem{arXiv2510} W.~Chen, Q.~Yang, S.~Shao, S.~Tang, Z.~Shi, and S.~Yu,
``Privacy-preserving semantic communication over wiretap channels with learnable
differential privacy,'' 2025, \emph{arXiv:2510.23274}.
\end{thebibliography}

\end{document}
